\section{Clause omission and polynomial concentration}\label{sec:concentration}

We now turn the integrated forcing estimate into a bound on the
fluctuations of the last satisfiable prefix. Omitting one clause can
delay the loss of satisfiability. The delay counts the prefixes where
that clause is pivotal, and its square counts pairs of pivotal times.
The key observation is that a future punctured prefix can survive only
if every intervening clause avoids killing the present solution set.

\subsection{Omitting a clause without changing time indices}

Fix an integer \(M\ge1\) and the independent clauses \(C_1,\ldots,C_M\).
For the next two lemmas, let
\[
 V=\max\{0\le m\le M:F_{n,m}\in\sat\}
\]
be the last satisfiable index capped at \(M\). For \(1\le i\le M\), set
\[
 V^{-i}=\max\left\{0\le m\le M:
       \bigwedge_{\substack{1\le j\le m\\j\ne i}}C_j\in\sat\right\}.
\]
The time index \(i\) is skipped; the remaining positions are not
renumbered. Thus \(V^{-i}\ge V\), and \(V^{-i}\) is independent of
\(C_i\).

The following coordinate-omission inequality belongs to the martingale
variance method associated with Efron and Stein~\cite{EfronStein} and
Steele's treatment of nonsymmetric statistics~\cite{Steele}.
Boucheron, Bousquet, Lugosi, and Massart
\cite[Section~3.2, equation~(3.6)]{BBLM} state the arbitrary-omission
form. The short proof fixes the deletion convention and constant.

\begin{lemma}[Deletion variance bound]\label{lem:variance-deletion}
For the capped indices just defined,
\[
 \Var(V)\le\sum_{i=1}^{M}\EE(V^{-i}-V)^2.
\]
\end{lemma}

\begin{proof}
Let \(\F_i=\sigma(C_1,\ldots,C_i)\), with \(\F_0\) trivial, and put
\[
 D_i=\EE[V\mid\F_i]-\EE[V\mid\F_{i-1}].
\]
Independence of the clauses and the fact that \(V^{-i}\) does not use
\(C_i\) give
\[
 \EE[V^{-i}\mid\F_i]=\EE[V^{-i}\mid\F_{i-1}].
\]
For \(Z_i=V-V^{-i}\), it follows that
\[
 D_i=\EE[Z_i\mid\F_i]-\EE[Z_i\mid\F_{i-1}].
\]
The tower property and conditional Jensen inequality yield
\[
 \EE D_i^2
 =\EE\bigl(\EE[Z_i\mid\F_i]\bigr)^2
  -\EE\bigl(\EE[Z_i\mid\F_{i-1}]\bigr)^2
 \le\EE Z_i^2.
\]
The martingale differences are orthogonal and sum to \(V-\EE V\).
Summing the displayed inequalities proves the lemma.
\end{proof}

\subsection{Pairs of pivotal times}

Fix \(i\). For \(i\le m\le M\), the punctured prefix and its killing
probability are
\[
 G_m=\bigwedge_{\substack{1\le j\le m\\j\ne i}}C_j,\qquad
 q_m=q_k(G_m).
\]
The function \(q_k\) was defined in \eqref{eq:kill}; in particular,
\(q_m=0\) when \(G_m\) is unsatisfiable.

\begin{lemma}[Pivotal pairs]\label{lem:pivotal}
For every \(1\le i\le M\),
\begin{equation}\label{eq:pivotal}
 \EE(V^{-i}-V)^2
 \le2\sum_{m=i}^{M}\EE\min\{Mq_m,1\}.
\end{equation}
\end{lemma}

\begin{proof}
Let
\[
 A_m=\{G_m\in\sat,\ G_m\wedge C_i\notin\sat\}.
\]
This event is exactly \(\{V<m\le V^{-i}\}\): prefix \(m\) is
satisfiable after omission and unsatisfiable before it. Hence, including
the endpoint \(m=M\),
\[
 V^{-i}-V=\sum_{m=i}^{M}\ind_{A_m}.
\]
For \(\ell\ge m\), on \(A_m\) the original prefix remains
unsatisfiable, so
\[
 A_m\cap A_\ell=A_m\cap\{G_\ell\in\sat\}.
\]

Condition on
\(\F_m^{(-i)}=\sigma(C_j:1\le j\le m,\ j\ne i)\).
The clause \(C_i\) and the future clauses
\(C_{m+1},\ldots,C_\ell\) are conditionally independent. The
conditional probability of \(A_m\) is \(q_m\). If \(G_m\) is
satisfiable, then \(G_\ell\) can be satisfiable only if none of these
future clauses individually makes \(G_m\) unsatisfiable. Each does so
with probability \(q_m\), independently. Thus
\begin{equation}\label{eq:pivotal-pair}
 \PP(A_m\cap A_\ell\mid\F_m^{(-i)})
 =q_m\PP(G_\ell\in\sat\mid\F_m^{(-i)})
 \le q_m(1-q_m)^{\ell-m}.
\end{equation}
If \(G_m\) is unsatisfiable, the relevant event and the right side are
both zero.

Expand the square of the indicator sum and bound its double sum by
twice the sum over \(m\le\ell\). For \(q>0\),
\[
 \sum_{\ell=m}^{M}(1-q)^{\ell-m}\le\min\{M,1/q\},
\]
because \(m\ge1\). Multiplication by \(q\), followed by expectation
and summation over \(m\), gives \eqref{eq:pivotal}. When \(q=0\), every
term in the corresponding product is zero, so no division by zero is
needed.
\end{proof}

\subsection{The variance bound}

The pivotal-pair estimate has reduced the squared omission delay to a
truncated killing probability. Taking a \(k\)-th root converts that
quantity into the forced fraction controlled by
Proposition~\ref{prop:forcing}.

\begin{proposition}[Polynomial concentration]\label{prop:concentration}
For every fixed integer \(k\ge3\), the statistic
\(V_n=\min(H_n-1,L_kn)\) from \eqref{eq:statistic} satisfies
\begin{equation}\label{eq:variance}
 \Var(V_n)=O_k(n^{1+2/k}).
\end{equation}
In particular, with \(\mu_n=\EE V_n/n\),
\begin{equation}\label{eq:concentration}
 \PP\bigl(|V_n-n\mu_n|\ge n^{1-\delta_k}\bigr)
 =O_k(n^{-\eta_k}),\qquad
 \delta_k=\frac{k-2}{4k},\quad \eta_k=\frac{k-2}{2k}.
\end{equation}
\end{proposition}

\begin{proof}
Set \(M=L_kn\), \(V=V_n\), and fix \(i\). With the notation of
Lemma~\ref{lem:pivotal}, let \(\rho_m=b(G_m)/n\).
The falling-factorial ratios in \eqref{eq:kill} give
\(q_m\le\rho_m^k\). Since \(\min\{z,1\}\le z^{1/k}\) for \(z\ge0\),
\[
 \min\{Mq_m,1\}\le(Mq_m)^{1/k}\le M^{1/k}\rho_m.
\]
Unconditionally, \(G_m\) has the law of \(F_{n,m-1}\).
Lemma~\ref{lem:pivotal} and Proposition~\ref{prop:forcing} therefore
give, uniformly in \(1\le i\le M\),
\[
 \EE(V^{-i}-V)^2
 \le2M^{1/k}\sum_{m=i}^{M}\EE\rho_m
 \le2M^{1/k}\sum_{j=0}^{M-1}\EE\frac{b(F_{n,j})}{n}
 =O_k(n^{2/k}).
\]
Summing over \(M=L_kn\) indices in
Lemma~\ref{lem:variance-deletion} proves \eqref{eq:variance}.
Chebyshev's inequality proves \eqref{eq:concentration}, because
\[
 1+\frac2k-2(1-\delta_k)
 =-\left(1-\frac2k-2\delta_k\right)=-\eta_k.
\]
\end{proof}

The exponent \(\delta_k\) is positive for \(k\ge3\), so the window
\(n^{1-\delta_k}\) is sublinear. The estimate is centered at the
finite-size mean \(n\mu_n\); comparing those means across sizes is a
separate task.
